% !TeX root = ../教材题库_全部条目_转写工作稿.tex
% Source: XJ-B1-C02-EX-P041-Y1378
% 第2章 / 2.1 相等关系与不等关系 / 教材 p.41、42 / PDF 第48、49页
% OCR未校对：公式与文字均可能有误；未完成数学TeX转写；禁止直接发给学生
% 原题及相邻原解截图（核对用）：../题目原图/XJ-B1-C02-EX-P041-Y1378-01.webp, ../题目原图/XJ-B1-C02-EX-P041-Y1378-02.webp
% 本文件只在转写完成后作为题干使用；OCR 条目尚不是可靠题干。
\begin{bankitem}{2.1 例题 例题9}{第2章 / 2.1 相等关系与不等关系 / 教材 p.41、42 / PDF 第48、49页}{XJ-B1-C02-EX-P041-Y1378}{OCR未校对：公式与文字均可能有误；未完成数学TeX转写；禁止直接发给学生}
\begin{ocrdraft}
[教材p.41 / PDF第48页]
\par
0某公司设计了如图2．1\_5所示的一块绿化
\par
景观地带，两条平行线段的两端用半圆弧相连接．已
\par
知这块绿化景观地带的内圈周长为4佣m,当平行线
\par
段的长设计为多少时，中间矩形区域的面积最大？
\par
解设平行线段长为工m，半圆的直径为dm，中
\par
图2．1一5
\par
间矩形区域的面积为Sm\ensuremath{\cdot}
\par
[教材p.42 / PDF第49页]
\par
由题意可知
\par
S=xd,
\par
且2十丌d一400，
\par
所以
\par
1
\par
1丌十2工
\par
20000
\par
2
\par
200
\par
当且仅当丌d=2=200，即=
\par
，一100时，等号成立．
\par
因此，当平行线段的长设计为100m时，中间矩形区域的面积S最大，最大值
\par
20000
\par
为
\end{ocrdraft}
\end{bankitem}
